\documentclass[10pt,a4paper]{amsart}
\usepackage{amsmath,amssymb,amsthm,graphicx}
\usepackage{verbatim}
\newenvironment{reference}{\comment}{\endcomment}
%\newenvironment{reference}{}{}
\newenvironment{slogan}{\comment}{\endcomment}
\newenvironment{history}{\comment}{\endcomment}
% For commutative diagrams we use Xy-pic
\usepackage[all]{xy}

% We use 2cell for 2-commutative diagrams.
\xyoption{2cell}
% Theorem environments.
%
\theoremstyle{plain}
\newtheorem{theorem}[subsection]{Theorem}
\newtheorem{proposition}[subsection]{Proposition}
\newtheorem{lemma}[subsection]{Lemma}

\theoremstyle{definition}
\newtheorem{definition}[subsection]{Definition}
\newtheorem{example}[subsection]{Example}
\newtheorem{exercise}[subsection]{Exercise}
\newtheorem{situation}[subsection]{Situation}

\theoremstyle{remark}
\newtheorem{remark}[subsection]{Remark}
\newtheorem{remarks}[subsection]{Remarks}

\numberwithin{equation}{subsection}

% Macros
%
\def\lim{\mathop{\mathrm{lim}}\nolimits}
\def\colim{\mathop{\mathrm{colim}}\nolimits}
\def\Spec{\mathop{\mathrm{Spec}}}
\def\Hom{\mathop{\mathrm{Hom}}\nolimits}
\def\Ext{\mathop{\mathrm{Ext}}\nolimits}
\def\SheafHom{\mathop{\mathcal{H}\!\mathit{om}}\nolimits}
\def\SheafExt{\mathop{\mathcal{E}\!\mathit{xt}}\nolimits}
\def\Sch{\mathit{Sch}}
\def\Mor{\mathop{\mathrm{Mor}}\nolimits}
\def\Ob{\mathop{\mathrm{Ob}}\nolimits}
\def\Sh{\mathop{\mathit{Sh}}\nolimits}
\def\NL{\mathop{N\!L}\nolimits}
\def\CH{\mathop{\mathrm{CH}}\nolimits}
\def\proetale{{pro\text{-}\acute{e}tale}}
\def\etale{{\acute{e}tale}}
\def\QCoh{\mathit{QCoh}}
\def\Ker{\mathop{\mathrm{Ker}}}
\def\Im{\mathop{\mathrm{Im}}}
\def\Coker{\mathop{\mathrm{Coker}}}
\def\Coim{\mathop{\mathrm{Coim}}}

% Boxtimes
%
\DeclareMathSymbol{\boxtimes}{\mathbin}{AMSa}{"02}

%
% Macros for moduli stacks/spaces
%
\def\QCohstack{\mathcal{QC}\!\mathit{oh}}
\def\Cohstack{\mathcal{C}\!\mathit{oh}}
\def\Spacesstack{\mathcal{S}\!\mathit{paces}}
\def\Quotfunctor{\mathrm{Quot}}
\def\Hilbfunctor{\mathrm{Hilb}}
\def\Curvesstack{\mathcal{C}\!\mathit{urves}}
\def\Polarizedstack{\mathcal{P}\!\mathit{olarized}}
\def\Complexesstack{\mathcal{C}\!\mathit{omplexes}}
% \Pic is the operator that assigns to X its picard group, usage \Pic(X)
% \Picardstack_{X/B} denotes the Picard stack of X over B
% \Picardfunctor_{X/B} denotes the Picard functor of X over B
\def\Pic{\mathop{\mathrm{Pic}}\nolimits}
\def\Picardstack{\mathcal{P}\!\mathit{ic}}
\def\Picardfunctor{\mathrm{Pic}}
\def\Deformationcategory{\mathcal{D}\!\mathit{ef}}
\usepackage{cleveref}
\crefname{originallemma}{original lemma}{original lemmas}
\Crefname{originallemma}{Original lemma}{Original lemmas}
\makeatletter
\newlabel{lemma-bound-cd}{{lemma-bound-cd}{source}{}{sourcepointer-1}{}}
\newlabel{lemma-bound-cd@cref}{{[originallemma][0][]lemma-bound-cd}{[1][0][]source}{}{}{}}
\makeatother
\crefname{originallemma}{original lemma}{original lemmas}
\Crefname{originallemma}{Original lemma}{Original lemmas}
\makeatletter
\newlabel{lemma-cd-change-rings}{{lemma-cd-change-rings}{source}{}{sourcepointer-2}{}}
\newlabel{lemma-cd-change-rings@cref}{{[originallemma][0][]lemma-cd-change-rings}{[1][0][]source}{}{}{}}
\makeatother
\crefname{originalreference}{original reference}{original references}
\Crefname{originalreference}{Original reference}{Original references}
\makeatletter
\newlabel{algebra-proposition-dimension}{{algebra-proposition-dimension}{source}{}{sourcepointer-3}{}}
\newlabel{algebra-proposition-dimension@cref}{{[originalreference][0][]algebra-proposition-dimension}{[1][0][]source}{}{}{}}
\makeatother
\crefname{originalreference}{original reference}{original references}
\Crefname{originalreference}{Original reference}{Original references}
\makeatletter
\newlabel{algebra-lemma-completion-Noetherian-Noetherian}{{algebra-lemma-completion-Noetherian-Noetherian}{source}{}{sourcepointer-4}{}}
\newlabel{algebra-lemma-completion-Noetherian-Noetherian@cref}{{[originalreference][0][]algebra-lemma-completion-Noetherian-Noetherian}{[1][0][]source}{}{}{}}
\makeatother
\crefname{originalreference}{original reference}{original references}
\Crefname{originalreference}{Original reference}{Original references}
\makeatletter
\newlabel{algebra-lemma-completion-complete}{{algebra-lemma-completion-complete}{source}{}{sourcepointer-5}{}}
\newlabel{algebra-lemma-completion-complete@cref}{{[originalreference][0][]algebra-lemma-completion-complete}{[1][0][]source}{}{}{}}
\makeatother
\crefname{originalreference}{original reference}{original references}
\Crefname{originalreference}{Original reference}{Original references}
\makeatletter
\newlabel{algebra-lemma-completion-faithfully-flat}{{algebra-lemma-completion-faithfully-flat}{source}{}{sourcepointer-6}{}}
\newlabel{algebra-lemma-completion-faithfully-flat@cref}{{[originalreference][0][]algebra-lemma-completion-faithfully-flat}{[1][0][]source}{}{}{}}
\makeatother
\crefname{originalreference}{original reference}{original references}
\Crefname{originalreference}{Original reference}{Original references}
\makeatletter
\newlabel{more-algebra-lemma-completion-dimension}{{more-algebra-lemma-completion-dimension}{source}{}{sourcepointer-7}{}}
\newlabel{more-algebra-lemma-completion-dimension@cref}{{[originalreference][0][]more-algebra-lemma-completion-dimension}{[1][0][]source}{}{}{}}
\makeatother
\crefname{originalreference}{original reference}{original references}
\Crefname{originalreference}{Original reference}{Original references}
\makeatletter
\newlabel{dualizing-lemma-ubiquity-dualizing}{{dualizing-lemma-ubiquity-dualizing}{source}{}{sourcepointer-8}{}}
\newlabel{dualizing-lemma-ubiquity-dualizing@cref}{{[originalreference][0][]dualizing-lemma-ubiquity-dualizing}{[1][0][]source}{}{}{}}
\makeatother
\crefname{originalreference}{original reference}{original references}
\Crefname{originalreference}{Original reference}{Original references}
\makeatletter
\newlabel{dualizing-lemma-special-case-local-duality}{{dualizing-lemma-special-case-local-duality}{source}{}{sourcepointer-9}{}}
\newlabel{dualizing-lemma-special-case-local-duality@cref}{{[originalreference][0][]dualizing-lemma-special-case-local-duality}{[1][0][]source}{}{}{}}
\makeatother
\crefname{originalreference}{original reference}{original references}
\Crefname{originalreference}{Original reference}{Original references}
\makeatletter
\newlabel{dualizing-lemma-nonvanishing-generically-local}{{dualizing-lemma-nonvanishing-generically-local}{source}{}{sourcepointer-10}{}}
\newlabel{dualizing-lemma-nonvanishing-generically-local@cref}{{[originalreference][0][]dualizing-lemma-nonvanishing-generically-local}{[1][0][]source}{}{}{}}
\makeatother
\title{The cohomological dimension of the maximal ideal of a local ring equals the dimension, and the top local cohomology of the ring is nonzero}
\author{The Stacks Project Authors}
\date{stacks-project:local-cohomology.tex:local-snapshot-20261012; GNU FDL 1.2 or later}
\begin{document}
\maketitle
\noindent\textit{Editorial note.} All delivered mathematical text is the author's own, copied from the pinned Stacks Project chapter source sources/stacks-local-cohomology/pinned-original-source.tex; no statement or proof step was rewritten, reordered or removed.\par\medskip
\noindent\textit{Editorial note.} Preamble provenance: the Stacks Project keeps one single shared preamble for every chapter; this card takes it directly from the staged sources/stacks-local-cohomology/preamble.tex as three verbatim spans (lines 8-12, 14-18 and 154-217), so no derived file was written and nothing was added or redefined. The wrapper class and the base packages are supplied by the assembler.\par\medskip
\noindent\textit{Editorial note.} Bibliography provenance: neither the counted statement nor its proof contains a citation (this lane rescans the whole span before writing the card), so the card carries no bibliography and no entry was reconstructed.\par\medskip
\noindent\textit{Editorial note.} Selection provenance: this card was opened under the file-based control contract (CONTROL.md), with the mechanical pre-card gate run on its statement and proof interval before the card was written.\par\medskip
\noindent\textit{Allowed standard prerequisites (editorial).}\begin{itemize}
\item The existence of an ideal of definition generated by d elements, where d is the dimension of the local ring, and the bound of the cohomological dimension by the number of generators. Both are invoked at the start of the counted proof (source lines 818-822), the second being a lemma of this chapter. Shows that the cohomological dimension is at most d, so that the remaining task is to exclude the strict inequality.
\item The behaviour of the cohomological dimension under the passage to the completion, and the existence of a normalized dualizing complex on a complete local ring with the local duality theorem. Invoked inside the counted proof (source lines 827-836 and 840-846) as pointers to lemmas of other chapters. Reduces the problem to a complete local ring, where local duality identifies the top local cohomology with a module that is known to be nonzero.
\end{itemize}\medskip
\ifdefined\hypertarget\hypertarget{sourcepointer-1}{}\fi
\noindent\textit{Source reference.} Original lemma lemma-bound-cd: sources/stacks-local-cohomology/pinned-original-source.tex, source lines 822--822. Navigation pointer only: the counted text refers to the Stacks item lemma-bound-cd, which belongs to another chapter of the project and is not part of this card. The pointer carries the target's own label as its number because this chapter's staged snapshot does not contain the other chapter's tag table; its content is not reproduced here.\par\medskip
\ifdefined\hypertarget\hypertarget{sourcepointer-2}{}\fi
\noindent\textit{Source reference.} Original lemma lemma-cd-change-rings: sources/stacks-local-cohomology/pinned-original-source.tex, source lines 835--835. Navigation pointer only: the counted text refers to the Stacks item lemma-cd-change-rings, which belongs to another chapter of the project and is not part of this card. The pointer carries the target's own label as its number because this chapter's staged snapshot does not contain the other chapter's tag table; its content is not reproduced here.\par\medskip
\ifdefined\hypertarget\hypertarget{sourcepointer-3}{}\fi
\noindent\textit{Source reference.} Original reference algebra-proposition-dimension: sources/stacks-local-cohomology/pinned-original-source.tex, source lines 820--820. Navigation pointer only: the counted text refers to the Stacks item algebra-proposition-dimension, which belongs to another chapter of the project and is not part of this card. The pointer carries the target's own label as its number because this chapter's staged snapshot does not contain the other chapter's tag table; its content is not reproduced here.\par\medskip
\ifdefined\hypertarget\hypertarget{sourcepointer-4}{}\fi
\noindent\textit{Source reference.} Original reference algebra-lemma-completion-Noetherian-Noetherian: sources/stacks-local-cohomology/pinned-original-source.tex, source lines 831--831. Navigation pointer only: the counted text refers to the Stacks item algebra-lemma-completion-Noetherian-Noetherian, which belongs to another chapter of the project and is not part of this card. The pointer carries the target's own label as its number because this chapter's staged snapshot does not contain the other chapter's tag table; its content is not reproduced here.\par\medskip
\ifdefined\hypertarget\hypertarget{sourcepointer-5}{}\fi
\noindent\textit{Source reference.} Original reference algebra-lemma-completion-complete: sources/stacks-local-cohomology/pinned-original-source.tex, source lines 832--832. Navigation pointer only: the counted text refers to the Stacks item algebra-lemma-completion-complete, which belongs to another chapter of the project and is not part of this card. The pointer carries the target's own label as its number because this chapter's staged snapshot does not contain the other chapter's tag table; its content is not reproduced here.\par\medskip
\ifdefined\hypertarget\hypertarget{sourcepointer-6}{}\fi
\noindent\textit{Source reference.} Original reference algebra-lemma-completion-faithfully-flat: sources/stacks-local-cohomology/pinned-original-source.tex, source lines 833--833. Navigation pointer only: the counted text refers to the Stacks item algebra-lemma-completion-faithfully-flat, which belongs to another chapter of the project and is not part of this card. The pointer carries the target's own label as its number because this chapter's staged snapshot does not contain the other chapter's tag table; its content is not reproduced here.\par\medskip
\ifdefined\hypertarget\hypertarget{sourcepointer-7}{}\fi
\noindent\textit{Source reference.} Original reference more-algebra-lemma-completion-dimension: sources/stacks-local-cohomology/pinned-original-source.tex, source lines 834--834. Navigation pointer only: the counted text refers to the Stacks item more-algebra-lemma-completion-dimension, which belongs to another chapter of the project and is not part of this card. The pointer carries the target's own label as its number because this chapter's staged snapshot does not contain the other chapter's tag table; its content is not reproduced here.\par\medskip
\ifdefined\hypertarget\hypertarget{sourcepointer-8}{}\fi
\noindent\textit{Source reference.} Original reference dualizing-lemma-ubiquity-dualizing: sources/stacks-local-cohomology/pinned-original-source.tex, source lines 842--842. Navigation pointer only: the counted text refers to the Stacks item dualizing-lemma-ubiquity-dualizing, which belongs to another chapter of the project and is not part of this card. The pointer carries the target's own label as its number because this chapter's staged snapshot does not contain the other chapter's tag table; its content is not reproduced here.\par\medskip
\ifdefined\hypertarget\hypertarget{sourcepointer-9}{}\fi
\noindent\textit{Source reference.} Original reference dualizing-lemma-special-case-local-duality: sources/stacks-local-cohomology/pinned-original-source.tex, source lines 844--844. Navigation pointer only: the counted text refers to the Stacks item dualizing-lemma-special-case-local-duality, which belongs to another chapter of the project and is not part of this card. The pointer carries the target's own label as its number because this chapter's staged snapshot does not contain the other chapter's tag table; its content is not reproduced here.\par\medskip
\ifdefined\hypertarget\hypertarget{sourcepointer-10}{}\fi
\noindent\textit{Source reference.} Original reference dualizing-lemma-nonvanishing-generically-local: sources/stacks-local-cohomology/pinned-original-source.tex, source lines 849--849. Navigation pointer only: the counted text refers to the Stacks item dualizing-lemma-nonvanishing-generically-local, which belongs to another chapter of the project and is not part of this card. The pointer carries the target's own label as its number because this chapter's staged snapshot does not contain the other chapter's tag table; its content is not reproduced here.\par\medskip
\begin{lemma}
\label{lemma-cd-maximal}
Let $(A, \mathfrak m)$ be a Noetherian local ring of dimension $d$.
Then $H^d_\mathfrak m(A)$ is nonzero and $\text{cd}(A, \mathfrak m) = d$.
\end{lemma}

\begin{proof}
By one of the characterizations of dimension, there exists
an ideal of definition for $A$ generated by $d$ elements, see
Algebra, Proposition \ref{algebra-proposition-dimension}.
Hence $\text{cd}(A, \mathfrak m) \leq d$ by
Lemma \ref{lemma-bound-cd}. Thus $H^d_\mathfrak m(A)$ is
nonzero if and only if $\text{cd}(A, \mathfrak m) = d$ if and only if
$\text{cd}(A, \mathfrak m) \geq d$.

\medskip\noindent
Let $A \to A^\wedge$ be the map from $A$ to its completion.
Observe that $A^\wedge$ is a Noetherian local ring of the
same dimension as $A$ with maximal ideal $\mathfrak m A^\wedge$.
See Algebra, Lemmas
\ref{algebra-lemma-completion-Noetherian-Noetherian},
\ref{algebra-lemma-completion-complete}, and
\ref{algebra-lemma-completion-faithfully-flat} and
More on Algebra, Lemma \ref{more-algebra-lemma-completion-dimension}.
By Lemma \ref{lemma-cd-change-rings}
it suffices to prove the lemma for $A^\wedge$.

\medskip\noindent
By the previous paragraph we may assume that $A$ is
a complete local ring. Then $A$ has a normalized dualizing complex
$\omega_A^\bullet$ (Dualizing Complexes, Lemma
\ref{dualizing-lemma-ubiquity-dualizing}).
The local duality theorem (in the form of
Dualizing Complexes, Lemma \ref{dualizing-lemma-special-case-local-duality})
tells us $H^d_\mathfrak m(A)$ is Matlis dual to
$\text{Ext}^{-d}(A, \omega_A^\bullet) = H^{-d}(\omega_A^\bullet)$
which is nonzero for example by
Dualizing Complexes, Lemma
\ref{dualizing-lemma-nonvanishing-generically-local}.
\end{proof}

\end{document}
